\chapter*{Appendix B: Prompts Used in the Pipeline}

The results in Chapter 5 depend heavily on how the models are asked, so the three prompts
that matter are printed here as they appear in the code.

\section*{B.1 Reading the whole board}

This prompt produced the change from 31.2 to 88.8 per cent board-content recall in Section
\ref{sec:board-reading}. The prompt it replaced asked the model for keywords visible on
the board.

\begingroup\footnotesize
\begin{verbatim}
This is a photograph of a classroom whiteboard during a university lecture.

Transcribe everything written on the whiteboard, exactly as written.

- Copy every number, name, identifier, date and symbol exactly. Do not round,
  correct or reformat values.
- Write tables as Markdown tables with the same rows and columns as the board.
- Write code, SQL and program text inside ``` fences, keeping line breaks.
- Write formulas as they appear. For a bar over a term write it as NOT(...),
  for example NOT(A+B).
- Keep diagrams short: name what they are and copy any labels on them, for
  example "Gate diagram: inputs A, B into a box labelled NAND, output NOT(AB)".
- If part of the board cannot be read, write [illegible]. Never guess, and never
  add anything that is not written on the board.
- Ignore any person in the picture. Describe only what is written.

Output only the transcription.
\end{verbatim}
\endgroup

\section*{B.2 Naming and transcribing each numbered box}

The Set-of-Mark prompt of Section 4.4.3. The boxes are drawn on the image by our own code
and numbered, and the model answers per number. Boxes it names \texttt{none} are dropped,
which is how a smudge or a sticker can be removed by the model even though the box finder
created it.

\begingroup\footnotesize
\begin{verbatim}
This is a classroom whiteboard from a university lecture. Numbered coloured boxes
have been drawn on it; each box surrounds one part of the writing. There are {n}
boxes, numbered 1 to {n}.

For every box, reply with:
- "id": the box number.
- "name": 1 to 4 words saying what kind of content the box holds, for example
  "Title", "Definition", "Truth table", "Gate symbol", "Block diagram",
  "SQL query", "Code", "List", "Worked example", "Formula". If the box holds no
  lecture writing (a smudge, a shadow, a logo or sticker, the edge of the board),
  use "none".
- "text": everything written inside that box, copied exactly: every number, name
  and symbol. For a bar over a term write NOT(...). Write a table as Markdown
  table rows. Write [illegible] for what cannot be read. Never guess.

Reply with JSON only: a list with one object per box, in box order, like
[{"id": 1, "name": "Title", "text": "Digital Logic Design"}]
\end{verbatim}
\endgroup

\section*{B.3 Writing one section of the notes}

These rules are shared by both languages. Rule 1 is the direct answer to the textbook
recitation described in Section 1.3, and rule 2 is what makes the notes point at the
numbered boxes.

\begingroup\footnotesize
\begin{verbatim}
RULES
1. Teach what THIS lecturer taught, with the lecturer's own examples and values.
   Explain it properly, for a student who missed the class: walk through each step
   on the board, say what it means and why the lecturer does it in that order, and
   spell out anything the board only shows as a symbol. What is not wanted is
   material from elsewhere: every fact in the explanation must come from the board
   or the transcript. Anything else belongs in the BACKGROUND section, never mixed
   in.
2. The board has numbered coloured boxes, listed below. Point the student at them,
   for example "Look at the orange box 3" or "the truth table in box 5". Use only
   the box numbers that are listed, and mention every listed box at least once.
3. Copy numbers, names, code, formulas and table values exactly as they appear in
   the box text. A wrong value is worse than a missing one.
4. The transcript is machine-made and has errors. Where it is garbled, rely on the
   board. Never invent what the lecturer said.
5. No greetings, no filler, no "in this section".
6. A table on the board is written as a Markdown table, never described row by row.
7. Plain text for formulas and symbols: A + B, (A + B)', NOT(A + B), A XOR B.
\end{verbatim}
\endgroup

The quotation instruction differs by language. The Banglish note quotes the lecturer's
real words, so the quotation can be checked against the transcript character by character,
and a quote that is not word for word in the transcript is deleted automatically. The
English note is read by a student who cannot read Banglish, so the lecturer is quoted in
English translation; a translation cannot be matched against a Banglish transcript, so the
automatic checker does not run on those files and each file records that fact in its
footer.

The example line in the Banglish quotation instruction, \texttt{"exact words from the
transcript"}, was itself copied into the notes as if it were a real quotation three times
across 26 files. The quotation checker could not catch it, because it looks for quotations
absent from the transcript and this is not a quotation at all, so a separate step now
removes the prompt's own placeholder text and counts how many it removed.
